\chapter{Controls and Commentary Engine}\label{ch:controls-engine}

\section{Role of the engine}

The controls engine is the declarative layer that turns the five
decision-point records of \cref{sec:decision-point-schema} into
executable rules. It sits between Stage 2 (anomaly detection) and
Stage 5 (FORTM narrative) of the AI pipeline. Its inputs are the
variance records, the GL journal extract, and the stakeholder
commentary responses; its outputs are routing actions and a
commentary register.

\section{Decision tree}

\Cref{fig:decision-tree} collapses the five gateways into a single
cascade. The tree is identical to the workflow in \cref{ch:workflow}
but omits task nodes; each leaf is a terminal action in the gateway
vocabulary.

\begin{figure}[htbp]
\centering
\begin{tikzpicture}[
  node distance=7mm and 8mm,
  gw/.style={diamond,draw,aspect=1.8,inner sep=1pt,align=center,
             font=\scriptsize,fill=tbamber!15,draw=tbamber!60,
             minimum height=9mm,minimum width=22mm},
  act/.style={draw,rounded corners=2pt,align=center,font=\scriptsize,
              fill=tbgreen!15,draw=tbgreen!60,minimum height=6mm},
  arr/.style={-{Latex[length=2mm]},thin}
]
  \node[gw] (dp1) {DP-001 \\ material?};
  \node[gw,below right=7mm and 4mm of dp1] (dp2) {DP-002 \\ journal?};
  \node[act,left=10mm of dp2] (noncrit) {no-commentary};
  \node[gw,below right=7mm and 4mm of dp2] (dp3) {DP-003 \\ investigate?};
  \node[act,left=8mm of dp3] (posted) {continue};
  \node[gw,below right=7mm and 4mm of dp3] (dp4) {DP-004 \\ M+4 received?};
  \node[act,left=8mm of dp4] (norinv) {continue};
  \node[gw,below right=7mm and 4mm of dp4] (dp5) {DP-005 \\ unexplained?};
  \node[act,left=8mm of dp5] (ontime) {continue};
  \node[act,below=of dp5] (summary) {prepare-summary};
  \node[act,right=20mm of dp5] (cat) {categorise};

  \draw[arr] (dp1) -- node[above,font=\tiny]{yes} (dp2);
  \draw[arr] (dp1) -- node[above,font=\tiny]{no}  (noncrit);
  \draw[arr] (dp2) -- node[above,font=\tiny]{not posted} (dp3);
  \draw[arr] (dp2) -- node[above,font=\tiny]{posted}     (posted);
  \draw[arr] (dp3) -- node[above,font=\tiny]{yes} (dp4);
  \draw[arr] (dp3) -- node[above,font=\tiny]{no}  (norinv);
  \draw[arr] (dp4) -- node[above,font=\tiny]{no}  (dp5);
  \draw[arr] (dp4) -- node[above,font=\tiny]{yes} (ontime);
  \draw[arr] (dp5) -- node[left,font=\tiny]{none}  (summary);
  \draw[arr] (dp5) -- node[above,font=\tiny]{some} (cat);
\end{tikzpicture}
\caption{Decision cascade encoded by the five decision-point
records.}\label{fig:decision-tree}
\end{figure}

\section{Rule format}

Each rule is a record with five fields:

\begin{lstlisting}[language=jsonx,caption={Generalised rule form.}]
{
  "id":          "<DP-nnn>",
  "state":       "<workflow stateId>",
  "input":       "<record field path>",
  "operator":    "<exceeds_threshold|boolean|in_set|...>",
  "threshold":   "<value or parameter name>",
  "branches":    [ {"condition":"...","action":"..."} ... ]
}
\end{lstlisting}

Rule evaluation is deterministic. When a record enters a gateway
state, the engine reads the named input field, applies the operator,
and dispatches the matching branch's action. Actions reference
downstream tasks in the workflow (\cref{ch:workflow}).

\section{Materiality thresholds}

\texttt{DP-001} reads \texttt{materiality\_limit} from the country
parameters sheet. During the AI rollout, the fitted ML threshold
(Stage 2) is compared side-by-side with the manual limit for six
months; divergences $> 10\%$ surface to the Head of Africa GFS for
review. Only after six clean cycles does the fitted threshold
replace the manual one in the parameter sheet.

\section{Journal posting}

\texttt{DP-002} has three branches. The engine first reads
\texttt{journal\_status} from the GL extract. If the status is
\textsc{posted}, the routing short-circuits to
\stateid{CHECK\_UNEXPLAINED}. If \textsc{posting\_required}, the
engine emits a posting request via
\stateid{SEND\_JOURNAL\_REQUEST} and returns to
\stateid{CHECK\_FURTHER\_INVESTIGATION}. If \textsc{not\_posted}
(stuck), the path proceeds directly to
\stateid{CHECK\_FURTHER\_INVESTIGATION}, skipping the request.

\section{Investigation and deadlines}

\texttt{DP-003} is a boolean gateway controlled by the
\texttt{investigation\_status} field. \texttt{DP-004} enforces the
\texttt{M+4} deadline by reading \texttt{days\_to\_deadline} from the
variance record. Crossing the deadline triggers the
\stateid{TRACK\_ENTRIES} branch and the risk-assessment update in
\stateid{UPDATE\_RISK}.

\section{Unexplained-variance handling}

\texttt{DP-005} is the final gate before summary. If
\texttt{unexplained\_variance\_count $= 0$}, the engine transitions to
\stateid{PREPARE\_SUMMARY}. Otherwise, the unexplained items are
passed to \stateid{CATEGORIZE\_UNEXPLAINED}, where each item is
tagged against the BSS risk vocabulary (\textsc{substantiated at risk},
\textsc{unsubstantiated unowned}, \textsc{unsubstantiated owned}) and
a commentary draft is requested.

\section{Commentary register}

All drafts produced by Stage 3 of the AI pipeline, together with the
analyst accept/edit decision, are persisted to the commentary
register. The register is the primary evidence artefact for
\texttt{EUC-002}: the register row count must equal the flagged
variance row count at the end of \stateid{UPDATE\_VARIANCE\_COMMENTARIES}.

\begin{table}[htbp]
\caption{Commentary register schema (derived from the variance
record).}\label{tab:commentary-register}
\centering
\small
\begin{tabularx}{\linewidth}{@{}llX@{}}
\toprule
\textbf{Field} & \textbf{Type} & \textbf{Notes} \\ \midrule
account\_id      & string  & GL account number \\
period           & string  & ISO yyyy-mm \\
variance\_amount & number  & Prior $\to$ current delta \\
draft            & string  & LLM-generated narrative \\
final            & string  & Post-review narrative \\
accept\_mode     & enum    & \textsc{as-is}, \textsc{edited}, \textsc{rejected} \\
reviewer         & string  & GFS Analyst name or role \\
review\_ts       & datetime & ISO-8601 \\
\bottomrule
\end{tabularx}
\end{table}

\section{Per-country parameter views}

The data product record carries a \texttt{x-country-views.HK} block
that injects an HKMA-aware vocabulary and sets
\texttt{defaultFilters.entity = "HKG"}. Adding a new country is a
matter of adding another view; no schema change is needed. Until
additional views exist, Hong Kong is the only LE with a fitted
threshold history.

\section{Observability}

Every rule firing is logged with input, branch taken, timestamp, and
the version of the ruleset. Logs are aggregated into the
\texttt{TB\_REVIEW\_VECTORS} table's metadata index so that prior
decisions are retrievable by account and period.
